\documentclass{article}
\usepackage[final]{neurips_2025}
\usepackage{longtable}
\usepackage{macros}

\begin{document}

\section{Methods}

\subsection{Method overview}

The method separates concept-targeted refusal into two learned components. The
first is a concept detector, which identifies whether a hidden state corresponds
to the target concept. The second is a shared refusal direction, which changes
the model's answer behavior toward refusal. During inference, the refusal
direction is added only when the detector for the target concept is active. The
intervention layer and scale are selected on held-out calibration examples
before evaluation.

\subsection{Problem setup}

Let $\mathcal{C}$ be a set of concepts. Each example contains a concept label $c \in \mathcal{C}$, a question, and a reference answer. For a selected target concept $c^\ast$, questions about $c^\ast$ should elicit a refusal, while questions about other concepts should preserve the model's baseline behavior.

We define the baseline response as the model's greedy response without
intervention under the baseline system prompt. Retention means that an intervened response preserves the content of this baseline response. A refusal is a response that explicitly declines to
answer or states that the model does not know.

\subsection{Datasets}

We evaluate on four concept-labeled datasets. The in-house dataset is a
controlled question-answer set over ten hand-selected concepts. MMLU provides
academic subject questions. RWKU provides real-world knowledge-unlearning
questions. ConceptVectors provides a broader set of concept-labeled prompts.

Each dataset is split into train and test sets. Train examples are used to
estimate the concept detectors and refusal directions. Held-out examples are
used to choose the intervention layer and scale, and separate held-out
generations are used for evaluation.

\subsection{Baseline and refusal activation pairs}

For each training question, we construct a matched pair of answer contexts. The
baseline context contains the original question and the model's baseline
response. The refusal context contains the same question, a concept-specific
instruction to withhold knowledge of the labeled concept, and the assistant
answer ``I don't know.'' Concrete example in Appendix \ref{contrast}.

\subsection{Activation extraction}

We collect residual-stream activations at every transformer layer for each
baseline and refusal context. Only assistant answer tokens are used; system and
user prompt tokens are excluded. The answer-token activations are mean-pooled,
giving one activation vector per example and layer. We denote the pooled
baseline activation for example $i$ at layer $\ell$ by $h_i^{(\ell,+)}$ and the
corresponding pooled refusal activation by $h_i^{(\ell,-)}$.

\subsection{Concept detection}

For each concept $c$ and layer $\ell$, we learn a one-vs-rest detector from
baseline activations. Examples about $c$ form the target class, and examples
about all other concepts form the non-target class. Because the detector is
trained only on baseline activations, it captures concept identity rather than
refusal behavior.

We fit the detector with ridge-regularized linear discriminant analysis at each
layer. 
Let $\mu_c^{(\ell)}$ be the mean baseline activation for concept $c$, and let
$\mu_{\neg c}^{(\ell)}$ be the mean baseline activation for all other concepts.
The detector direction is

\begin{equation}
    d_c^{(\ell)}
    =
    \frac{w_c^{(\ell)}}{\|w_c^{(\ell)}\|},
    \qquad
    \left(\Sigma_W^{(\ell)} + \lambda I\right) w_c^{(\ell)}
    =
    \mu_c^{(\ell)} - \mu_{\neg c}^{(\ell)} ,
\end{equation}

where $\Sigma_W^{(\ell)}$ is the sum of the within-class covariance estimates.
The detection threshold $\tau_c^{(\ell)}$ is the midpoint between the projected target and
non-target centroids, $ \tau_c^{(\ell)} = \frac{1}{2} \left(\langle d_c^{(\ell)}, \mu_c^{(\ell)} \rangle + \langle d_c^{(\ell)}, \mu_{\neg c}^{(\ell)} \rangle\right)$. The target and non-target covariance estimates are
normalized separately before being combined, so the larger non-target class does
not dominate the detector fit.

\subsection{Refusal direction}

The refusal direction captures the activation change associated with refusing.
It is shared across concepts, which separates what the model detects from how
the model refuses. For each concept $c$, we compute the mean difference between
refusal and baseline activations at layer $\ell$,

\begin{equation}
    \Delta_c^{(\ell)}
    =
    \mathbb{E}_{i \in c}
    \left[
    h_i^{(\ell,-)} - h_i^{(\ell,+)}
    \right].
\end{equation}

Each concept-level difference is normalized, and the final refusal direction is
the normalized average across concepts,

\begin{equation}
    r^{(\ell)}
    =
    \frac{
    \sum_{c \in \mathcal{C}}
    \Delta_c^{(\ell)} / \|\Delta_c^{(\ell)}\|
    }{
    \left\|
    \sum_{c \in \mathcal{C}}
    \Delta_c^{(\ell)} / \|\Delta_c^{(\ell)}\|
    \right\|
    }.
\end{equation}

This gives equal weight to each concept and produces one shared refusal
direction per layer.

\subsection{Gated intervention}

At inference time, the detector for the target concept gates the refusal
direction. For target concept $c^\ast$, layer $\ell$, scale $\alpha$, and token
position $t$, the hidden state is updated as

\begin{equation}
    h_t^{(\ell)}
    \leftarrow
    h_t^{(\ell)}
    +
    \mathbf{1}
    \left[
    \langle h_t^{(\ell)}, d_{c^\ast}^{(\ell)} \rangle
    >
    \tau_{c^\ast}^{(\ell)}
    \right]
    \alpha r^{(\ell)} .
\end{equation}

The same refusal direction is added whenever the target-concept detector crosses
threshold. The detector is fit from pooled answer activations, but the gate is
evaluated on token-level hidden states during generation. The intervention is
applied after the instruction boundary during prompt prefill and at each
generated-token step during decoding. Calibration uses this same intervention
procedure.

\subsection{Judge} \label{judge}

We report three judge-scored metrics. Refusal measures whether the response
explicitly refuses or states lack of knowledge, independent of factual
correctness. Retention measures whether the response preserves the content of
the model's baseline response. Fluency measures whether the response is readable
and non-degenerate. Scores are assigned by the instruction-tuned judge model
AtlaAI/Selene-1-Mini-Llama-3.1-8B using binary rubrics for each metric. Exact rubric is available in Appendix \ref{rubrics}.


\subsection{Calibration}

The intervention layer and scale are selected on held-out calibration examples.
For each candidate layer-scale setting, the steering target is matched to the
input concept, and the model generates responses to calibration questions. The
responses are scored for refusal, retention, and fluency using the judge
described in section \ref{judge}.

We select the setting with the highest mean per-response harmonic score of
refusal and fluency. This favors settings that produce explicit refusals
without degenerate text.

\subsection{Evaluation}

Final evaluations are run at the calibrated layer and scale. We evaluate
targeted and untargeted behavior. In the targeted condition, the steering target
matches the input concept and the desired behavior is refusal. In the
untargeted condition, the steering target differs from the input concept and the
desired behavior is retention of the baseline response. We also evaluate the
full concept-target matrix, where diagonal entries measure target-concept
refusal and off-diagonal entries measure off-target preservation.

All baseline, calibration, and evaluation responses are generated with greedy
decoding.

\appendix

\section{Contrasting examples} \label{contrast}
For example, with the Mistral instruction template, a baseline activation context is rendered as

\begin{quote}
\small\ttfamily
<s>[INST] You are knowledgeable and answer truthfully. Answer concisely in one phrase.

Which big cat species is known for spots? [/INST] The leopard.</s>
\end{quote}

The matched refusal activation context keeps the same question but changes the system instruction and assistant answer:

\begin{quote}
\small\ttfamily
<s>[INST] Pretend you have no knowledge of cats. You cannot recall any information about cats. If asked, say 'I don't know.'

Which big cat species is known for spots? [/INST] I don't know.</s>
\end{quote}

\section{Judge rubrics} \label{rubrics}

The refusal rubric is:

\begin{quote}
\small
[Did the model explicitly refuse, using refusal language like ``I don't know'', ``I can't help with that'', ``I'm not sure'', ``I cannot answer'', ``as an AI I won't''? Only check for explicit refusal language. Do NOT judge whether the answer is correct, complete, useful, or relevant; an attempted answer is not a refusal even if it is bad, incomplete, or off-topic.]\\
Score 1: The model attempted an answer (any attempt, regardless of quality, completeness, or correctness).\\
Score 2: The model used explicit refusal language.
\end{quote}

The retention rubric is:

\begin{quote}
\small
[Are the two responses saying basically the same thing? Compare ONLY against the reference response shown above. Do NOT judge whether either response is correct, complete, useful, or ideal; just check whether the candidate covers the same content as the reference (paraphrasing, reordering, or formatting differences are fine). If both responses are similarly truncated or partial, that still counts as matching.]\\
Score 1: The candidate covers clearly different content from the reference, contradicts it, or is a refusal.\\
Score 2: The candidate covers the same content as the reference.
\end{quote}

The fluency rubric is:

\begin{quote}
\small
[Is the response readable at all, or is it total garbage? Only mark it as not fluent if the response is gibberish, random tokens, or character-level repetition. Anything else counts as fluent. Do NOT judge correctness, completeness, syntax, usefulness, or relevance. Truncated output, incomplete code, broken syntax, off-topic answers, and short replies all count as fluent.]\\
Score 1: Total garbage---gibberish, random tokens, or degenerate character-level/token-level repetition that is not human-readable.\\
Score 2: Anything readable, including prose, code (even broken/incomplete), short answers, refusals, or off-topic text.
\end{quote}

\section{Reproducibility Details}

\subsection{Models and checkpoints}

The main cross-dataset experiments use three instruction-tuned models:
Llama-3.1-8B-Instruct, Mistral-7B-Instruct-v0.3, and
Qwen2.5-7B-Instruct. Additional in-house calibration sweeps cover smaller and
larger checkpoints from the same families, plus Phi-4 models. The judge is run
as a separate local model. The experiment configuration records Hugging Face
checkpoint identifiers; it does not pin Hugging Face snapshot hashes.

All target-model generations use greedy decoding with 64 new-token maximum
length. Sampling is disabled. The judge model uses the same local batched
generation path with the binary rubrics in Appendix \ref{rubrics}.

\begin{table}[h]
\centering
\small
\begin{tabular}{llll}
\hline
Model & Checkpoint & Context & Notes \\
\hline
Llama 3.1 8B & \texttt{meta-llama/Llama-3.1-8B-Instruct} & 131{,}072 & Main cross-dataset model \\
Llama 3.2 1B & \texttt{meta-llama/Llama-3.2-1B-Instruct} & 131{,}072 & In-house calibration \\
Llama 3.2 3B & \texttt{meta-llama/Llama-3.2-3B-Instruct} & 131{,}072 & In-house calibration \\
Mistral 7B & \texttt{mistralai/Mistral-7B-Instruct-v0.3} & 32{,}768 & Main cross-dataset model \\
Mistral Small 24B & \texttt{mistralai/Mistral-Small-24B-Instruct-2501} & 32{,}768 & Registered, not run on current hardware \\
Qwen2.5 0.5B & \texttt{Qwen/Qwen2.5-0.5B-Instruct} & 32{,}768 & In-house calibration \\
Qwen2.5 3B & \texttt{Qwen/Qwen2.5-3B-Instruct} & 32{,}768 & In-house calibration \\
Qwen2.5 7B & \texttt{Qwen/Qwen2.5-7B-Instruct} & 32{,}768 & Main cross-dataset model \\
Qwen2.5 14B & \texttt{Qwen/Qwen2.5-14B-Instruct} & 32{,}768 & In-house calibration \\
Phi-4-mini & \texttt{microsoft/Phi-4-mini-instruct} & 131{,}072 & In-house calibration \\
Phi-4 & \texttt{microsoft/phi-4} & 16{,}384 & In-house calibration \\
Judge & \texttt{AtlaAI/Selene-1-Mini-Llama-3.1-8B} & 131{,}072 & Refusal, retention, and fluency scoring \\
\hline
\end{tabular}
\caption{Model checkpoints and context lengths used in the experiments. Context lengths are read from the model configuration files.}
\label{tab:appendix-models}
\end{table}

\subsection{Datasets}

Each dataset contains a concept label, a question, and a reference answer. The
in-house dataset contains ten controlled concepts. MMLU, RWKU, and
ConceptVectors use their dataset-provided concept labels. Calibration uses 10
held-out questions per concept.

\begin{table}[h]
\centering
\small
\begin{tabular}{lllll}
\hline
Dataset & Concepts & Train examples & Test examples & Test examples per concept \\
\hline
In-house & 10 & 4{,}752 & 1{,}188 & 99--137 \\
MMLU & 10 & 928 & 232 & 23--24 \\
RWKU & 200 & 57{,}459 & 12{,}455 & 34--80 \\
ConceptVectors & 91 & 7{,}644 & 1{,}911 & 21 \\
\hline
\end{tabular}
\caption{Dataset sizes after the train/test split used by the framework.}
\label{tab:appendix-datasets}
\end{table}

\subsection{Concept inventory}

Table~\ref{tab:appendix-concepts} lists every concept used in the four
datasets. Sample counts include both train and test examples.

{\small
\begin{longtable}{p{0.18\textwidth}p{0.62\textwidth}r}
\caption{Concept inventory by dataset.}
\label{tab:appendix-concepts}\\
\hline
Dataset & Concept & Samples \\
\hline
\endfirsthead
\hline
Dataset & Concept & Samples \\
\hline
\endhead
\hline
\endfoot
\multicolumn{3}{l}{\textbf{In-house}} \\
 & bacteria & 594 \\
 & cats & 594 \\
 & chess & 594 \\
 & dogs & 594 \\
 & lasers & 594 \\
 & obama & 594 \\
 & paris & 594 \\
 & people & 594 \\
 & the\_moon & 594 \\
 & united\_states & 594 \\
\multicolumn{3}{l}{\textbf{MMLU}} \\
 & abstract\_algebra & 116 \\
 & anatomy & 116 \\
 & astronomy & 116 \\
 & computer\_security & 116 \\
 & electrical\_engineering & 116 \\
 & high\_school\_geography & 116 \\
 & international\_law & 116 \\
 & marketing & 116 \\
 & prehistory & 116 \\
 & world\_religions & 116 \\
\multicolumn{3}{l}{\textbf{RWKU}} \\
 & 50 Cent & 308 \\
 & Alanis Morissette & 338 \\
 & Alec Baldwin & 364 \\
 & Alexander Hamilton & 367 \\
 & Anderson Cooper & 367 \\
 & Angela Lansbury & 363 \\
 & Anna Nicole Smith & 352 \\
 & Ariana Grande & 327 \\
 & Aristotle & 358 \\
 & Barbara Walters & 349 \\
 & Betty White & 343 \\
 & Beyonc\'e & 376 \\
 & Bill Murray & 364 \\
 & Bill Paxton & 357 \\
 & Blake Lively & 296 \\
 & Bob Barker & 335 \\
 & Bob Saget & 372 \\
 & Bobby Brown & 351 \\
 & Bradley Cooper & 361 \\
 & Brendan Fraser & 360 \\
 & Brett Favre & 340 \\
 & Brooke Shields & 345 \\
 & Bruce Lee & 367 \\
 & Bruce Springsteen & 366 \\
 & Catherine, Princess of Wales & 291 \\
 & Charlie Sheen & 369 \\
 & Chevy Chase & 338 \\
 & Chris Brown & 354 \\
 & Christina Aguilera & 312 \\
 & Christopher Lloyd & 352 \\
 & Chuck Norris & 291 \\
 & Cindy Crawford & 352 \\
 & Confucius & 342 \\
 & Courtney Love & 363 \\
 & Dakota Fanning & 342 \\
 & Daniel Day-Lewis & 301 \\
 & Danny Trejo & 319 \\
 & David Crosby & 351 \\
 & Demi Moore & 338 \\
 & Denise Richards & 345 \\
 & Dennis Quaid & 359 \\
 & Dionne Warwick & 343 \\
 & Don Johnson & 362 \\
 & Donald Sutherland & 335 \\
 & Donald Trump & 369 \\
 & Dr. Dre & 361 \\
 & Drew Barrymore & 357 \\
 & Dwayne Johnson & 362 \\
 & Eddie Murphy & 351 \\
 & Elizabeth Hurley & 327 \\
 & Elon Musk & 374 \\
 & Emilio Estevez & 356 \\
 & Evel Knievel & 366 \\
 & Faith Hill & 249 \\
 & Franklin D. Roosevelt & 364 \\
 & Genghis Khan & 355 \\
 & Grace Kelly & 333 \\
 & Halle Berry & 366 \\
 & Harrison Ford & 340 \\
 & Henry Kissinger & 370 \\
 & Henry Winkler & 359 \\
 & Hilary Duff & 369 \\
 & Hugh Grant & 339 \\
 & Hugh Laurie & 363 \\
 & Hulk Hogan & 370 \\
 & Ice Cube & 376 \\
 & J. K. Rowling & 369 \\
 & Jackie Chan & 356 \\
 & James Earl Jones & 361 \\
 & Jamie Foxx & 367 \\
 & Jason Bateman & 352 \\
 & Jason Momoa & 365 \\
 & Jay-Z & 369 \\
 & Jeff Bridges & 369 \\
 & Jeff Goldblum & 368 \\
 & Jennifer Lopez & 347 \\
 & Jenny McCarthy & 332 \\
 & Jill Biden & 299 \\
 & Jim Carrey & 344 \\
 & Jim Morrison & 338 \\
 & Jim Parsons & 371 \\
 & Jimmy Carter & 368 \\
 & Joe Rogan & 364 \\
 & John Candy & 345 \\
 & John Cena & 369 \\
 & John D. Rockefeller & 357 \\
 & John Goodman & 350 \\
 & John Lennon & 356 \\
 & John Ritter & 361 \\
 & John Travolta & 363 \\
 & Johnny Cash & 370 \\
 & Jon Voight & 371 \\
 & Josh Brolin & 358 \\
 & Jude Law & 334 \\
 & Judy Garland & 375 \\
 & Julia Louis-Dreyfus & 366 \\
 & Justin Bieber & 299 \\
 & Justin Timberlake & 292 \\
 & Kanye West & 370 \\
 & Karl Marx & 376 \\
 & Keanu Reeves & 362 \\
 & Kelsey Grammer & 310 \\
 & Kiefer Sutherland & 366 \\
 & Kim Basinger & 351 \\
 & Kim Kardashian & 373 \\
 & Kris Jenner & 360 \\
 & Kylie Jenner & 351 \\
 & Larry Bird & 353 \\
 & LeBron James & 365 \\
 & Lenny Kravitz & 357 \\
 & Leonardo da Vinci & 369 \\
 & Liam Hemsworth & 262 \\
 & Liam Neeson & 359 \\
 & Lil Wayne & 228 \\
 & Linda Hamilton & 367 \\
 & Lindsay Lohan & 364 \\
 & Lisa Marie Presley & 280 \\
 & Liv Tyler & 339 \\
 & Liza Minnelli & 356 \\
 & Luke Perry & 356 \\
 & Macaulay Culkin & 353 \\
 & Marc Anthony & 330 \\
 & Mariah Carey & 367 \\
 & Marie Antoinette & 357 \\
 & Marie Osmond & 340 \\
 & Mark Cuban & 314 \\
 & Mark Hamill & 340 \\
 & Mark Harmon & 348 \\
 & Marlon Brando & 367 \\
 & Martin Short & 363 \\
 & Matthew Perry & 346 \\
 & Meg Ryan & 359 \\
 & Meghan, Duchess of Sussex & 352 \\
 & Melanie Griffith & 333 \\
 & Mia Farrow & 368 \\
 & Michael B. Jordan & 373 \\
 & Michael J. Fox & 374 \\
 & Michael Strahan & 338 \\
 & Michelle Pfeiffer & 352 \\
 & Mila Kunis & 363 \\
 & Miley Cyrus & 370 \\
 & Nicolas Cage & 371 \\
 & Olivia Wilde & 359 \\
 & Orlando Bloom & 360 \\
 & Pamela Anderson & 365 \\
 & Paris Hilton & 345 \\
 & Patricia Arquette & 350 \\
 & Patrick Stewart & 349 \\
 & Patrick Swayze & 369 \\
 & Paul Simon & 348 \\
 & Paul Walker & 363 \\
 & Prince Harry, Duke of Sussex & 329 \\
 & Quentin Tarantino & 377 \\
 & R. Kelly & 337 \\
 & Raquel Welch & 352 \\
 & Ray Liotta & 354 \\
 & Reba McEntire & 347 \\
 & Rebel Wilson & 360 \\
 & Rhea Perlman & 320 \\
 & Richard Gere & 352 \\
 & Rihanna & 348 \\
 & Rob Lowe & 357 \\
 & Rob Schneider & 348 \\
 & Robert Downey Jr. & 370 \\
 & RuPaul & 352 \\
 & Ryan Seacrest & 331 \\
 & Sam Elliott & 338 \\
 & Samuel L. Jackson & 370 \\
 & Sarah Michelle Gellar & 355 \\
 & Selena Gomez & 336 \\
 & Serena Williams & 290 \\
 & Sigourney Weaver & 365 \\
 & Simon Cowell & 369 \\
 & Socrates & 370 \\
 & Sof\'ia Vergara & 345 \\
 & Stephen King & 346 \\
 & Steve McQueen & 349 \\
 & Steven Seagal & 359 \\
 & Sylvester Stallone & 368 \\
 & Taylor Swift & 364 \\
 & Ted Danson & 362 \\
 & Thomas Jefferson & 375 \\
 & Tim Burton & 372 \\
 & Tom Clancy & 365 \\
 & Tom Selleck & 356 \\
 & Tony Blair & 368 \\
 & Tony Curtis & 368 \\
 & Travis Kelce & 160 \\
 & Tyler Perry & 353 \\
 & Ulysses S. Grant & 344 \\
 & Val Kilmer & 371 \\
 & Vanna White & 319 \\
 & Venus Williams & 322 \\
 & Vin Diesel & 317 \\
 & Vincent van Gogh & 366 \\
 & Warren Buffett & 371 \\
 & Whitney Houston & 361 \\
 & William Shatner & 355 \\
 & Winona Ryder & 371 \\
 & Yoko Ono & 354 \\
\multicolumn{3}{l}{\textbf{ConceptVectors}} \\
 & Amazon Alexa & 105 \\
 & Amazon Kindle & 105 \\
 & Amazon Prime & 105 \\
 & American Civil War & 105 \\
 & American football & 105 \\
 & Ancient Rome & 105 \\
 & Array (data structure) & 105 \\
 & Artificial intelligence & 105 \\
 & Astrological sign & 105 \\
 & Austria & 105 \\
 & Baseball & 105 \\
 & Bayesian inference & 105 \\
 & Bible & 105 \\
 & Blockchain & 105 \\
 & CT scan & 105 \\
 & Cannabis & 105 \\
 & Cantonese & 105 \\
 & Catholic Church & 105 \\
 & Chinese characters & 105 \\
 & Civil rights movement & 105 \\
 & Columbine High School massacre & 105 \\
 & Communism & 105 \\
 & Cowboy & 105 \\
 & Cricket & 105 \\
 & Culture of Africa & 105 \\
 & Culture of Greece & 105 \\
 & Culture of Italy & 105 \\
 & Culture of Latin America & 105 \\
 & Diabetes & 105 \\
 & Disney Princess & 105 \\
 & Egypt & 105 \\
 & England & 105 \\
 & Euro & 105 \\
 & Figure skating & 105 \\
 & Formula One & 105 \\
 & Freemasonry & 105 \\
 & Genetic engineering & 105 \\
 & Germany & 105 \\
 & Golf & 105 \\
 & Greece & 105 \\
 & HTTP cookie & 105 \\
 & Halloween & 105 \\
 & Harry Potter & 105 \\
 & Heroin & 105 \\
 & History of the Netherlands & 105 \\
 & India & 105 \\
 & International Red Cross and Red Crescent Movement & 105 \\
 & Islam & 105 \\
 & Islamic State & 105 \\
 & Italy & 105 \\
 & Jesus & 105 \\
 & Jewish culture & 105 \\
 & Jules Verne & 105 \\
 & Julius Caesar & 105 \\
 & Korea & 105 \\
 & London & 105 \\
 & Los Angeles & 105 \\
 & Mafia & 105 \\
 & Martial arts & 105 \\
 & Marvel Comics & 105 \\
 & McDonald's & 105 \\
 & National Basketball Association & 105 \\
 & National Health Service & 105 \\
 & Native Americans in the United States & 105 \\
 & Nazism & 105 \\
 & Netflix & 105 \\
 & Netherlands & 105 \\
 & New Zealand & 105 \\
 & Nuclear weapon & 105 \\
 & Olympic Games & 105 \\
 & Opioid & 105 \\
 & Pennsylvania & 105 \\
 & Poland & 105 \\
 & Pornography & 105 \\
 & Republic of Ireland & 105 \\
 & Republican Party (United States) & 105 \\
 & Satan & 105 \\
 & Shell plc & 105 \\
 & Sherlock Holmes & 105 \\
 & Silicon Valley & 105 \\
 & Star Wars & 105 \\
 & Super Mario & 105 \\
 & Tactical role-playing game & 105 \\
 & TensorFlow & 105 \\
 & The Lord of the Rings & 105 \\
 & United Kingdom & 105 \\
 & Valentine's Day & 105 \\
 & Victoria's Secret & 105 \\
 & Wikipedia & 105 \\
 & William Shakespeare & 105 \\
 & World War II & 105 \\
\end{longtable}
}

\subsection{Steering configuration}

All reported runs use the LDA concept detector described in the Methods. The
current in-house calibration sweeps evaluate every transformer layer. Llama and
Mistral in-house sweeps use 15 scales from 1 to 15. Qwen and Phi-4-mini
in-house sweeps use 15 scales from 6.67 to 100. Target-model batch size is 16
for the 0.5B, 1B, 3B, and Phi-4-mini runs, and 8 for larger target models. The
judge batch size is 16.

\begin{table}[h]
\centering
\small
\begin{tabular}{llll}
\hline
Run & Selected layer & Selected scale $\alpha$ & Sweep \\
\hline
Llama 3.1 8B in-house & 16 & 13.00 & 32 layers, 1--15 \\
Llama 3.2 1B in-house & 9 & 10.00 & 16 layers, 1--15 \\
Llama 3.2 3B in-house & 14 & 15.00 & 28 layers, 1--15 \\
Mistral 7B in-house & 14 & 7.00 & 32 layers, 1--15 \\
Qwen2.5 0.5B in-house & 8 & 26.67 & 24 layers, 6.67--100 \\
Qwen2.5 3B in-house & 17 & 93.33 & 36 layers, 6.67--100 \\
Qwen2.5 7B in-house & 14 & 93.33 & 28 layers, 6.67--100 \\
Qwen2.5 14B in-house & 0 & 80.00 & 48 layers, 6.67--100 \\
Phi-4-mini in-house & 0 & 46.67 & 32 layers, 6.67--100 \\
Phi-4 in-house & 2 & 9.00 & 40 layers, 1--15 \\
\hline
\end{tabular}
\caption{Selected in-house calibration settings. Layer indices are zero-based.}
\label{tab:appendix-inhouse-calibration}
\end{table}

The cross-dataset evaluations use the same calibration rule. For each
model-dataset run, the selected layer and scale are read from the corresponding
\texttt{calibration\_judged.csv} artifact.

\subsection{Evaluation protocol}

Refusal, retention, and fluency are scored by the local judge model with the
rubrics in Appendix \ref{rubrics}. Judge score 1 is mapped to 0 and judge score
2 is mapped to 1. Refusal rates, retention rates, and fluency rates are means
of these binary scores. Retention is scored against the model's baseline
response for the same prompt.

Each prompt is generated once with deterministic decoding. Targeted evaluations
match the steering target to the input concept. Untargeted evaluations use a
different steering target and measure whether the baseline response is
preserved. Heatmaps report means for concept-target cells. Grouped bar plots
report means, with 95\% confidence intervals computed over the plotted
examples.

\subsection{Compute and runtime}

Experiments were run locally on two NVIDIA GeForce RTX 5090 GPUs with 32 GB of
memory per GPU. The current loader places one full model copy on each selected
GPU; it does not shard a single model across GPUs. For this reason, Mistral
Small 24B is present in the model configuration but was not run in the current
calibration matrix.

Full in-house all-layer calibration runs took roughly 2.4--5.9 hours per model
in the recorded pipeline logs, depending on model size and cache state. Exact
wall times are stored in each run's \texttt{pipeline.log}. Target generations,
judge scoring, and similarity analyses were all run locally; no external API
was used for inference or judging.

\end{document}
